% The triple-tail latch's second and third stages as flow/comp/subckt.py builds them (its first
% stage is the double tail's with NMOS cascodes): cross-coupled PMOS on a clocked PMOS tail (twice
% latch_on), pulled down by NMOS on the first stage's drains P and Q (rst); then cross-coupled
% inverters on a clocked NMOS tail (latch_on), held high by PMOS on the second stage's nodes (four
% times rst).
% The preamble of every isometric drawing: standalone where installed (as docs/images/*.tex),
% else article with the page cropped to the drawing. Then Latin Modern sans and the blocks.
\IfFileExists{standalone.cls}{\documentclass[tikz,border=2mm]{standalone}\newenvironment{drawing}{}{}}{%
  \documentclass{article}\usepackage{tikz}\pagestyle{empty}\newsavebox\drawingbox
  \newenvironment{drawing}{
  \begin{lrbox}{\drawingbox}}{
    \end{lrbox}%
    \pagewidth=\dimexpr\wd\drawingbox+4mm\relax\pageheight=\dimexpr\ht\drawingbox+\dp\drawingbox+4mm\relax
\hoffset=-1in\voffset=-1in\shipout\vbox{\kern2mm\hbox{\kern2mm\usebox\drawingbox}}}}
\usepackage[T1]{fontenc}
\usepackage{lmodern}
\renewcommand{\familydefault}{\sfdefault}
\usepackage{thesis-drawings}

\usepackage{thesis-circuits}
\begin{document}
\begin{drawing}
  \begin{tikzpicture}[circuit]
    % stage 2: pull-downs on P and Q, cross-coupled PMOS, clocked PMOS tail
    \draw (-1.6,1.60) node[nmos](d1){};
    \draw (1.6,1.60) node[nmos, xscale=-1](d2){};
    \draw (d1.S) node[ground]{} (d2.S) node[ground]{};
    \draw (d1.G) node[left, sig]{$P$};
    \draw (d2.G) node[right, sig]{$Q$};
    \draw (-1.6,3.50) node[pmos, xscale=-1](u1){};
    \draw (1.6,3.50) node[pmos](u2){};
    \draw (d1.D) -- (u1.D) (d2.D) -- (u2.D);
    \coordinate (Mp) at ($(d1.D)!0.5!(u1.D)$); \coordinate (Mn) at ($(d2.D)!0.5!(u2.D)$);
    \draw (u1.G) -- (Mn) (u2.G) -- (Mp);
    \fill (Mp) circle (1.6pt) (Mn) circle (1.6pt);
    \draw (-1.6,4.27) -- (1.6,4.27);
    \draw (0,5.04) node[pmos](mt2){};
    \draw (mt2.D) -- (0,4.27);
    \fill (0,4.27) circle (1.6pt);
    \draw (mt2.S) node[vcc]{};
    \draw (mt2.G) node[left, sig]{$\overline{\mathrm{CLK}}$};
    \noteR{mt2}{$2\times$on}
    \noteR{d1}{rst}
    \noteL{u1}{on}
    % stage 3: clocked NMOS tail, cross-coupled inverters, held high by PMOS on the second stage
    \draw (0,8.10) node[nmos](mt3){};
    \draw (mt3.S) node[ground]{};
    \draw (mt3.G) node[left, sig]{CLK};
    \noteR{mt3}{on}
    \draw (mt3.D) -- (0,8.87);
    \draw (-1.6,8.87) -- (1.6,8.87);
    \fill (0,8.87) circle (1.6pt);
    \draw (-1.6,9.64) node[nmos, xscale=-1](n1){};
    \draw (1.6,9.64) node[nmos](n2){};
    \draw (-1.6,11.70) node[pmos, xscale=-1](p1){};
    \draw (1.6,11.70) node[pmos](p2){};
    \draw (p1.S) node[vcc]{} (p2.S) node[vcc]{};
    \draw (n1.D) -- (p1.D) (n2.D) -- (p2.D);
    \draw (n1.G) -- (p1.G) (n2.G) -- (p2.G);
    \coordinate (X) at ($(n1.D)!0.7!(p1.D)$);
    \coordinate (Y) at ($(n2.D)!0.3!(p2.D)$);
    \draw (X) -- (X -| n1.G) node[jump crossing]{} -- (X -| n2.G);
    \draw (Y) -- (Y -| n2.G) node[jump crossing]{} -- (Y -| n1.G);
    \fill (X) circle (1.6pt) (Y) circle (1.6pt) (X -| n2.G) circle (1.6pt) (Y -| n1.G) circle (1.6pt);
    \noteL{n1}{on}
    \noteL{p1}{init}
    \draw (-3.4,11.30) node[pmos](h1){};
    \draw (3.4,11.30) node[pmos, xscale=-1](h2){};
    \draw (h1.S) node[vcc]{} (h2.S) node[vcc]{};
    \draw (h1.D) |- (X) (h2.D) |- (Y);
    \draw (h1.G) -| ($(h1.G) - (0.4,0)$) |- (Mp);
    \draw (h2.G) -| ($(h2.G) + (0.4,0)$) |- (Mn);
    \noteR{h1}{$4\times$rst}
    \draw (X) node[above left, sig]{$X$};
    \draw (Y) node[below right, sig]{$Y$};
  \end{tikzpicture}
\end{drawing}
\end{document}
